% =====================================================================
\chapter{从意图到可验证执行}\label{ch:intent}
% =====================================================================

本章处理问题 \ref{prob:trust}：如果编码成本趋近于零、人不再逐行审阅实现，正确性如何
被机械地保证。答案由三部分构成：把正确拆成可分别处理的维度，把中间过程定义为一条
可验证细化链，把信任组织成两种生成模式的复合结构。本章同时给出该链条上一个不可机械
封闭的入口，以及它在何种意义上可以被共识逼近。

\section{正确性的维度}\label{sec:five-correctness}

一个 1\,ns 的调用被 AI 改写成 1\,$\mu$s 的调用，功能完全一致但性能差了一千倍，这个
例子说明"正确"不是单一维度。

\begin{definition}[五个正确性维度]\label{def:five-correctness}
见表 \ref{tab:five-correctness}。
\end{definition}

\begin{table}[htbp]
\centering
\small
\caption{五个正确性维度及其可机械检查性}\label{tab:five-correctness}
\begin{tabular}{L{2.3cm}L{6.6cm}L{5.5cm}}
\toprule
\textbf{维度} & \textbf{需要回答的问题} & \textbf{能否机械检查} \\
\midrule
功能正确性 & 输出与状态变化是否符合要求？ & 对明确规格，可证明或测试 \\
副作用正确性 & 有没有执行未授权的 I/O、写入、网络请求、状态变更？ &
  可通过效果规格、权限边界与运行时隔离限制 \\
资源正确性 & 时间、内存、带宽、能耗是否满足预算？ &
  部分可静态证明；部分只能测量或在指定环境下给出统计保证 \\
系统正确性 & 组件组合后，并发、故障恢复、状态一致性是否仍正确？ &
  部分可形式化证明，部分需模型检查、故障注入与运行验证 \\
目标表达正确性 & 机器执行的规格是否真的是人想要的目标？ &
  \textbf{不能从程序正确性证明中自动推出} \\
\bottomrule
\end{tabular}
\end{table}

表 \ref{tab:five-correctness} 最后一行是整个问题的关键：\emph{一份程序完全可能正确地
实现了一份错误的规格}。编译器即使百分之百正确，也可以把一个错误的程序百分之百正确
地编译成机器码；形式化证明同样无法拯救一个遗漏关键要求的规格。这一区分的历史表述
是软件需求工程中的验证与确认之分\cite{boehm1984}。

\section{可验证细化链}\label{sec:chain}

\begin{definition}[可验证细化链]\label{def:chain}
\term{可验证细化链}由七个阶段构成：
\begin{enumerate}[label=(G\arabic*)]
\item \term{目的与意图}：要实现什么；哪些结果是成功；哪些副作用不可接受；目标冲突时
      如何取舍。
\item \term{系统与环境模型}：把运行中的状态、事件、操作、资源、操作系统接口、外部
      环境与可观察行为表示出来。
\item \term{形式化行为契约}：按定义 \ref{def:contract} 给出允许行为、禁止行为、
      必须成立的不变量、进展条件、资源预算与环境假设。
\item \term{可实现性判定}：检查契约是否自相矛盾，目标在指定资源、硬件、操作系统与
      环境条件下是否真正可实现。
\item \term{实现综合与搜索}：由生成器、搜索算法或其他综合器产生候选实现。候选实现
      不预先获得信任。
\item \term{从抽象行为到二进制的正确性证明}：检查候选实现及其编译、链接与其他转换，
      证明生成的机器码在规定环境下满足行为契约。
\item \term{经认证的部署与实际执行}：核实正在运行的是已验证的二进制，在规定环境中
      执行，并对无法静态保证的条件进行监测、限制与故障处理。
\end{enumerate}
\end{definition}

这条链的意义不在于增加七个流程步骤，而在于\emph{给每一步规定不同的数学任务}。
(G2) 不应是自然语言描述，而是对系统状态、行为与环境建立可解释的模型；(G3) 判断要求
是否可满足；(G4) 搜索可能的实现；(G5) 建立从实现到所要求行为的证明；(G6) 确保实际
执行的对象与证明所对应的对象一致。

\begin{principle}[边界上传递什么]\label{pr:boundary}
在每一道边界上，不仅传递一个产物，还要传递\emph{这个产物代表什么}，以及
\emph{它为什么可以被接受的证据}。
\end{principle}

\begin{remark}[不可机械替代的入口]
验证结果与反例可以反向进入 (G2)、(G3)、(G4)；唯有 (G1) 中的意图审定，无法被一般性的
代码正确性证明取代。理由见第 \ref{sec:intent-gap} 节。
\end{remark}

\section{转换边与证据}\label{sec:five-edges}

把定义 \ref{def:chain} 的七个阶段压成五条转换边，每条边有其典型证据与关键缺口。
表 \ref{tab:edges} 给出。

\begin{table}[htbp]
\centering
\footnotesize
\caption{五条转换边、典型证据与关键缺口}\label{tab:edges}
\begin{tabular}{L{1.5cm}L{3.4cm}L{4.6cm}L{4.9cm}}
\toprule
\textbf{边} & \textbf{输入 $\to$ 输出} & \textbf{典型证明/证据} & \textbf{关键假设与主要缺口} \\
\midrule
I$\to$R & 自然语言/对话 $\to$ 结构化需求、领域事实与待澄清项 &
  利益相关者审定、需求来源链接、例子与反例、冲突报告 &
  意图没有自动给定的数学真值；外部事实与利益冲突不能只靠类型检查解决 \\
R$\to$S & 需求 $\to$ 形式化性质/契约 &
  追溯矩阵、需求级测试、一致/可满足/可实现分析、真空性与变异分析 &
  形式化映射要同时保留\emph{禁例与正例}；需求遗漏无法只从内部一致性推出 \\
S$\to$P & 规格 $\to$ 程序/系统候选 &
  综合算法、SMT、模型检查、Hoare/类型/依赖类型证明、反例引导归纳综合 &
  有限 DSL、逻辑片段与搜索预算；必须防止前提过强与规格空洞 \\
P$\to$C$\to$B & 源程序 $\to$ 逐阶段中间表示 $\to$ 机器码/可部署制品 &
  经验证编译器定理、逐次翻译验证/认证、证明携带代码 &
  必须定义源/目标语义；链接、装载、汇编、运行时、第三方库与硬件未必在证明范围内 \\
B$\to$E & 二进制与配置 $\to$ 真实环境行为 &
  部署完整性、系统验证、运行时监控、隔离与能力机制、性能与调度分析 &
  模型与真实 CPU、操作系统、网络、调度、故障与服务假设必须对齐 \\
\bottomrule
\end{tabular}
\end{table}

\begin{remark}[关系 $R_i$ 的异质性]
按注 \ref{rmk:source}，不同转换边上的解释关系不必是同一类。对 I$\to$R 与 R$\to$S
而言，可用的"关系"可能只是追溯映射而非形式化关系\cite{mucha2024,mokos2020}；
对 P$\to$C$\to$B 而言，它可能是编译器的语义保持定理\cite{leroy2009,cakeml2019}、
单次翻译验证\cite{necula2000}或翻译认证\cite{krijnen2024}；
对 B$\to$E 而言，它可能是运行时监控的判定\cite{taleb2023}。把它们强行统一为
同一种关系是错误的做法。
\end{remark}

\begin{remark}[每条边上已有的工程实现]
表 \ref{tab:edges} 并非纯粹的设想：P$\to$C$\to$B 上的经验证编译
\cite{leroy2009}与经验证编译器后端\cite{cakeml2019}、内核级的精化证明
\cite{klein2009}、以及证明携带代码\cite{necula1997}已经在特定语言、编译器、
内核与系统中取得了重要进展。真正未解决的是：如何把这些高可信度的局部成果，
扩展到更广泛的目标、更多种类的运行环境，以及实现可以被任意生成或动态调整的系统。
\end{remark}

\section{组合条件与失效模式}\label{sec:composition}

命题 \ref{prop:soundness} 已经给出组合的充分条件。这里补充三个具体的组合失效反例，
它们的作用是证明"局部正确"确实不蕴含"整体正确"。

\begin{example}[三个组合失效]\label{ex:compose-fail}
\begin{enumerate}[label=(\alph*)]
\item 源程序语义正确，但链接器选错了库。
\item 每个服务本身满足接口契约，但契约对网络重试的假设互不兼容。
\item 功能证明成立，但时延预算没有覆盖缓存争用。
\end{enumerate}
三者的共同结构是：命题 \ref{prop:soundness} 的条件 (a) 或 (b) 之一被破坏。
\end{example}

用契约代数（定义 \ref{def:operad}、命题 \ref{prop:operad-compose}）读这三个反例，
它们不再彼此无关，而是三类失败：

\begin{corollary}[三个反例的共同结构]\label{cor:three-failures}
\begin{enumerate}[label=(\alph*)]
\item 链接器选错库：\emph{着色失配}，被链接实现其颜色不等于链接点声明的颜色；
\item 契约对网络重试的假设互不兼容：\emph{接口图出现环}，重试把无环调用图变成
      带环图，而着色未随环重新验证（注 \ref{rmk:feedback}）；
\item 时延预算未覆盖缓存争用：\emph{契约分量缺失}，资源分量 $R$ 未进入颜色，
      代价被当作实现细节。
\end{enumerate}
三者分别是颜色外延不足、接口图改变、契约分量缺失。
\end{corollary}

结论是：组合规则必须把这些接口条件与\emph{非功能性条件}显式纳入。

\begin{remark}[保证的整体论性质]\label{rmk:duhem}
推论 \ref{cor:three-failures}(c) 与注 \ref{rmk:monitor-limit} 共同说明：任何验证结论
实际断言的不是"性质 $\varphi$ 成立"，而是合取式
\[
\varphi \;\land\; A_1 \;\land\; A_2 \;\land \cdots \land A_k,
\]
其中 $A_i$ 为全部环境假设。$k$ 事先不可枚举，只有当某个 $A_i$ 真被违反时它才显形。
这是 Duhem--Quine 论题\cite{quine1951}在软件验证中的形态：不存在对单一命题的判决性
检验，只有对整个合取式的检验。它的直接后果是，"测试通过"的强度严格受限于
\emph{未被枚举的假设集}。这也解释了表 \ref{tab:evidence} 为何必须为每类证据附一份
"不应声称"清单，该清单就是已知的 $A_i$。
\end{remark}

\section{意图鸿沟}\label{sec:intent-gap}

这是本章最需要谨慎的一节。

\begin{problem}[意图鸿沟]
人最初用自然语言表达的意图 $I$，如何确保被转化为正确的形式化规格 $\Phi$？我们希望
有 $I \equiv \Phi$，这里的等价指"$\Phi$ 所要求的行为完整地反映了人的真实目标"。
\end{problem}

一个流行的论证是：由 Gödel 不完备定理与 Rice 定理，意图与规格之间的裂缝是
\emph{原则性}的、\emph{永远无法机械封死}的。该论证不成立。

\begin{proposition}[本文不接受该论证的理由]\label{prop:no-godel}
Gödel 不完备定理讨论的是满足特定条件的形式系统中\emph{形式可证明性}的限制，它并不
直接说明机器无法判断一份需求规格是否符合某个人的真实意图。Rice 定理
\cite{rice1953}说明对一般程序而言非平凡语义性质不可判定，这是对通用静态分析的
根本限制，但它不是关于人的心理意图与形式规格之间关系的定理。要从中推出"意图永远
无法完整形式化"，还需要额外的哲学与认知假设。Turing 的原始结果\cite{turing1936}
同样只适用于明确编码的计算问题与一般算法主张。
\end{proposition}

那么意图鸿沟存在吗？存在，但根据是\emph{认识论的}，不是\emph{逻辑的}。

\begin{proposition}[意图鸿沟的三条认识论根据]\label{prop:three-grounds}
\begin{enumerate}[label=(\roman*)]
\item \term{需求是生成性的，而非枚举性的}。人在看到系统行为之前无法完整列出自己的
      需求。"不能泄露隐私"这条要求，在看到系统确实在收集数据之前可能根本没有浮现。
      这不是逻辑不可判定，而是人类认知的构造性局限：\emph{偏好是被环境激活的，
      不是预先存在的完整清单}。
\item \term{规格是意图的有损压缩}。意图存在于一个因果的、具身的、演化塑造的心智中；
      规格是一个有限的符号对象。有损压缩意味着信息必然丢失，而丢掉的是哪些部分
      事先不可知。这是信息论层面的现象，不是元数学层面的。
\item \term{验证与确认的区分至今未被推翻}。验证（verification）回答"实现是否符合
      规格"；确认（validation）回答"系统是否满足真实需求"。这个区分之所以存在并
      沿用至今，正是因为后者无法被机械地归约为前者。
\end{enumerate}
\end{proposition}

把命题 \ref{prop:no-godel} 与 \ref{prop:three-grounds} 合起来，得到本文对该问题的
准确表述：

\begin{principle}[意图鸿沟的准确表述]\label{pr:gap}
规格转换过程不是一个可以\emph{普遍}机械保证的步骤，原因不在逻辑不可判定，而在人的
偏好是动态的、情境性的、在遭遇之前不可全知的。
\end{principle}

这一修正没有削弱命题，只是替换了承重墙。它同时说明：机器可以辅助澄清需求、提出反例、
生成场景、寻找矛盾，并检查规格内部是否一致；而最终采用哪份规格、它是否完整表达了
目标，仍然需要一个明确的责任主体。

\begin{remark}[三个问题，而不是一个；以及自动化能为每一个做什么]\label{rmk:three-questions}
意图鸿沟常被当作单一障碍来讨论。它其实是三个，它们需要不同的工具，而把它们分开来看
就会清楚：本文框架拒绝机械化的区域，比通常以为的要小。

\begin{enumerate}[label=(\alph*)]
\item \term{语义表达}：给定事实上已经确定下来的需求，它们能否被写成形式化规格？在有界
      领域中这大体上可以自动化，其工具是类型系统、定理证明器，以及第 \ref{ch:literature}
      章讨论的规格综合方法。命题 \ref{prop:three-grounds} 没有任何一条否认这一点，把它读成
      否认这一点是过度解读。
\item \term{需求发现}：利益相关者是否已经找到并表达了所有要紧的东西？这正是命题
      \ref{prop:three-grounds}(i)，其补救手段是交互：反例、场景，以及第
      \ref{ch:comprehension} 章的探测循环。它部分可自动化，意义是机器可以生成探针；
      它不可自动化，意义是答案必须来自有切身利害的人。
\item \term{规范性的取舍}：当需求彼此冲突——时延与安全性、隐私与分析能力、成本与
      性能——谁有资格作出决定？这根本不是认识论问题；它是第 \ref{sec:authority} 节的
      权威轴，其工具是程序性的：一个被指定的责任机构，以及一份决定记录。
\end{enumerate}
这样读这一鸿沟，也纠正了相反方向上一个常见的过度表述。因为 (a) 大体可机械化，
"某些意图可以被精确定义并自动检查"是真的；不能被机械化的是残留部分——尚未被激活的
偏好、尚未浮现的冲突，以及无人有资格独自作出的决定。因此本文所捍卫的准确主张比
"意图无法被形式化"要窄：\emph{一份关于实现正确性的形式化证明，本身无法确立该规格是否
服务于所有相关各方的目标}，而开放世界中的剩余部分必须由一个独立的确认过程来处理。
该过程可以使用自动化工具；它不得把一个工具的内部一致性当作"目标是对的"的证据。
\end{remark}

\begin{remark}[意图鸿沟是框架问题]\label{rmk:frame}
命题 \ref{prop:three-grounds} 的三条根据不是三项独立发现，它们共同指向一个更古老的
问题。要形式化意图，必须预先判定哪些事实\emph{相关}；而判定相关性需要已经知道哪些
结论有待推出，这正是框架问题：不预先知道答案，就无法枚举相关前提。

因此意图鸿沟不是形式化方法的缺陷，而是框架问题的一个实例。这一识别有两项后果：
\begin{enumerate}[label=(\alph*)]
\item 它解释了为什么更强的形式化工具不解决该问题，工具提高的是\emph{推导}能力，
      而问题出在\emph{前提的枚举}；
\item 它解释了为什么\emph{交互}是唯一有效的手段。框架问题在一般情况下未解，只在具备
      感知--行动闭环的系统中被回避：智能体无需预先枚举相关性，因为环境会对动作给出
      反馈。第 \ref{ch:comprehension} 章的探测循环与本节的可追溯精化，都是把这一
      回避机制搬到需求层。
\end{enumerate}
于是"交互式精化"是被逼出来的，而不是在若干候选方案中被选中的。
\end{remark}

\begin{remark}[无视集是框架问题被逐层级偿付之处]\label{rmk:frame-per-level}
注 \ref{rmk:frame} 把意图鸿沟与框架问题等同起来，就到此为止。这一等同还可以再推进一步，
而推进之后就把本文此前彼此分离的两部分接了起来。

本文所使用的框架问题，其形式是\emph{判定相关性}的困难：在还不知道要推出什么结论之前，
就要说出哪些事实与结论有关。无视集（定义 \ref{def:oblivious}）正是一个这种形状的
相关性判断，只不过它被显式写出、并且是局部的：$\Obl(\Lambda_i)$ 命名那些 $\Lambda_i$
声明与自己无关的差异。因此两者之间的关系不是类比，而是尺度上的关系——而把"得到了什么"
与"没有得到什么"说清楚是值得的。

得到的第一样东西是\emph{局部化}。为一个系统整体提出的框架问题，不会因为改成对它的
某个层级提出而变小；但它被该层级所声明的观察映射约束住，并且用它能够陈述的"货币"来
清偿。全局问题——哪些事实与系统的目的相关——仍然没有一般解，本文也不声称有。层级结构
换来的是：这个全局问题永远不会被一次性提出；每个层级回答自己有界的那一版本，并把
答案公布出来。

得到的第二样东西是\emph{可检验性}。仅仅被持有的相关性判断无法被批评；被写下来的
则可以。第 \ref{sec:neardecomp} 节给出一个已声明的判断可能出错的两种方式——接口对外部
行为可能不充分，或者观察映射相对于事实上可观察的东西可能不完备。两者都是检查，而不是
意见。在这个意义上，无视集是框架问题的工程形态：它不是框架问题的解，而是一个为它
付出代价、且这笔代价可以被审计的地方。

由此得到两项后果，它们也正是本条注放在这里、而不放在第 \ref{ch:comprehension} 章的
原因。第一，它解释了原则 \ref{pr:descent} 的下潜协议为什么有那样一条停止规则：一直
下潜到支配该变更的不变量变得可见为止，因为那个深度正是相关的无视集已被咨询过的深度，
而继续下潜所换来的相关性，没有任何义务依赖其上。第二，它使意图层的困难可以用与模块
边界相同的措辞读出来：意图鸿沟是 $\Lk{\infty}$ 上的框架问题，而 $\Lk{\infty}$ 恰恰是
不能为机器声明任何无视集的层级，因为声明一个无视集本身就已经是在决定目的。
\end{remark}

\section{验证与生成的不对称}\label{sec:asymmetry}

另一个需要修正的表述是把"判卷比答卷容易"称为"$\mathrm{P}$ 与 $\mathrm{NP}$ 的世界性
红利"。这一说法不成立：$\mathrm{P}$ 是否等于 $\mathrm{NP}$ 仍是开放问题，我们并不知道
一般而言找到解是否必然比检查解困难。

但必须把这一修正与另一件事区分开来：$\mathrm{P}$ 与 $\mathrm{NP}$ 的未决，
\emph{并不}使"验证比生成便宜"这一现象成为虚构。

\begin{proposition}[结构性不对称]\label{prop:asym-struct}
多项式时间证书检查的结构性不对称是 $\mathrm{NP}$ 类的\emph{定义本身}：一个解可以被
高效验证，这件事本身是关于计算的一类基本事实，它不依赖于我们是否证明了
$\mathrm{P} \neq \mathrm{NP}$。
\end{proposition}

\begin{proof}[论证]
$\mathrm{NP}$ 的定义为"存在多项式时间的验证器 $V$ 与多项式长度的证书 $w$，使得
$x \in L \iff \exists w: V(x,w) = 1$"。这一定义本身断言了验证的存在性与高效性，
而不对寻找 $w$ 的复杂度作任何断言。因此验证器架构的有效性建立在定义层面，与
$\mathrm{P}$ 对 $\mathrm{NP}$ 问题的最终答案无关。
\end{proof}

这一成本量级提示了"验证便宜"的准确边界。限制不在 $\mathrm{P}$ 与
$\mathrm{NP}$ 的未决，而是\emph{证书长度}：检查一份证明可以在多项式时间内完成，
而这份证明本身可以长到天文数字。因此：

\begin{principle}[验证成本的双分量]\label{pr:two-costs}
"验证比生成便宜"必须被分解为两个独立命题：
\begin{enumerate}[label=(\alph*)]
\item \term{检查时间便宜}：给定证书，判定是高效的（命题 \ref{prop:asym-struct}）；
\item \term{证书尺寸便宜}：证书本身足够短，可以被存储、传输与复核。
\end{enumerate}
(a) 由 $\mathrm{NP}$ 的定义保证；(b) \emph{不}由它保证，且在许多实际系统中是真正的
瓶颈。一个只满足 (a) 的验证体系在工程上不可用，因为无人能复核一份无法读完的证书。
\end{principle}

原则 \ref{pr:two-costs}(b) 是可测量的：对给定的证明器与性质类，测量证书尺寸的分布
即可。它同时解释了为什么"小内核 + 大证明"的架构在实践中总要配合一个\emph{可读的}
证明摘要层（不变量、契约的散文陈述、反例），而不是只提供证明项。

\begin{example}[形式化验证的历史成本量级]\label{ex:sel4-cost}
微内核 seL4 的初始功能正确性证明投入了约 11 人年（不计可复用工具与框架的投资），
当时的证明脚本约 20 万行，用于验证约 8500 行 C 代码\cite{klein2009}。这个量级说明
形式化验证的历史成本结构。AI 可能显著降低编写证明、寻找不变量与修复证明的成本，
但规格设计、模型选择、证明架构与可信基础仍然是实际工作。
\end{example}

\section{证据等级}\label{sec:evidence-levels}

不同的证据提供不同强度的支持，混淆它们是这一领域最常见的错误。
表 \ref{tab:evidence} 给出证据等级与相应的"不应声称"清单。

\begin{table}[htbp]
\centering
\footnotesize
\caption{证据等级}\label{tab:evidence}
\begin{tabular}{L{2.6cm}L{5.6cm}L{6.2cm}}
\toprule
\textbf{证据类型} & \textbf{能够支持的结论} & \textbf{不应声称的结论} \\
\midrule
机器检查的形式证明 & 对声明的命题，在给定逻辑、定义、公理、库及检查器假设下成立 &
  不证明命题完整表达意图；不覆盖未建模环境；不自动说明时延或能耗 \\
经验证编译器定理/翻译证书 & 特定源/目标语义间存在所定义的保持或精化关系，或单次
  翻译结果经验证 & 不意味着链接器、操作系统、硬件与外部服务都经过验证 \\
静态分析与抽象解释 & 相对于可靠近似，证明或保守估计所分析的程序性质 &
  不保证精确、无误报；不等于实际性能测量 \\
测试与变异测试 & 发现反例、评估已列测试上的规格行为，提供有限样本证据 &
  不证明所有输入与所有执行安全；通过率不是规格语义等价的通用证明 \\
统计性能证据 & 指定数据、采样流程、硬件与置信方法下的经验分布 &
  不等于硬实时最坏执行时间；不能无条件外推到分布变化后的部署 \\
运行时监控 & 对被观察到的轨迹及监控器语义给出在线判定 &
  若观察不完整，不得把"未检测到"表述为"所有执行均满足" \\
\bottomrule
\end{tabular}
\end{table}

\begin{principle}[三个必须区分的状态]\label{pr:three-states}
"检测到违反"、"尚未观察到违反"与"所有执行均被证明安全"是三个不同的状态，绝不能
互相替代。
\end{principle}

\begin{definition}[制品证据清单]\label{def:evidence-package}
每个制品应携带十项证据，见表 \ref{tab:evidence-package}。该清单的作用是使"某个二进制
是否可被信任"成为一个可审计的问题，而不是一个关于作者身份的问题。
\end{definition}

\begin{table}[htbp]
\centering
\small
\caption{制品证据清单}\label{tab:evidence-package}
\begin{tabular}{L{0.8cm}L{13.4cm}}
\toprule
\textbf{编号} & \textbf{证据项} \\
\midrule
1 & 规格版本与需求追溯标识 \\
2 & 形式化规格与环境假设 \\
3 & 候选源代码哈希 \\
4 & 编译器、参数、目标架构与依赖版本 \\
5 & 证明及检查器版本 \\
6 & 资源分析模型与结果 \\
7 & 测试与监控策略 \\
8 & 可信计算基清单 \\
9 & \textbf{已知未覆盖行为} \\
10 & \textbf{重新验证规则}（哪些变化触发重新计算受影响的证据）\\
\bottomrule
\end{tabular}
\end{table}

第 9 项与第 10 项是最容易缺失、也最重要的两项：前者把"保证的边界"显式化，后者把
保证从一次性检查变成生命周期机制。

\section{副作用与约束机制}\label{sec:effects}

问题 \ref{prob:trust} 中的例子混合了不同种类的副作用。分开看才知道各自的机械机制。

\begin{definition}[副作用四分类]\label{def:effects}
见表 \ref{tab:effects}。
\end{definition}

\begin{table}[htbp]
\centering
\small
\caption{副作用分类与相应的机械机制}\label{tab:effects}
\begin{tabular}{L{1.8cm}L{4.6cm}L{5.0cm}L{3.0cm}}
\toprule
\textbf{类别} & \textbf{例子} & \textbf{机械机制} & \textbf{验证时点} \\
\midrule
功能性 & 生成的函数修改了不该修改的全局状态 &
  类型系统、效果类型（把"会碰什么"变成签名的一部分）、所有权系统 &
  编译期，零运行时成本 \\
资源性 & 1\,ns 变 1\,$\mu$s、内存暴涨 &
  燃料/计价计量、最坏执行时间分析、配额 &
  编译期 $+$ 运行时双重 \\
安全性 & 越权访问、逃逸沙箱 &
  能力模型、沙箱、字节码验证器 &
  准入时 $+$ 运行时 \\
语义漂移 & 规格本身就错了 &
  \textbf{没有机械机制} & 人的审定与共识 \\
\bottomrule
\end{tabular}
\end{table}

前三类已有机械机制覆盖，且是纵深配置（证明 $>$ 类型 $>$ 测试 $>$ 计价 $>$
隔离 $>$ 监控，每一层的信任基础都比上一层更简单）。只有第四类（规格偏离意图）永远
需要人。

\begin{remark}[三类工程原型的保证范围]\label{rmk:three-prototypes}
三种常见工程机制覆盖的是不同种类的性质，不能并列宣传：
\begin{enumerate}[label=(\alph*)]
\item 浏览器沙箱保证的是\emph{执行隔离与权限边界}，不保证功能正确性。它是地球上
      最大的不可信代码运行时：每天执行数以亿计来源不明、从未审查的代码，它从不审查，
      它隔离。
\item 字节码验证器验证的是\emph{其规则所覆盖的安全性质}（无越界访问、循环有界、
      无悬挂指针），不是任意业务逻辑正确性。一个几万行的确定性验证器裁决着百万行自由
      代码的准入，这正是"编译器式的、不以意志为转移的保证"的实例。
\item 资源计价机制保证的是\emph{抽象资源计量下的约束}，不是真实的墙上时钟时间。
      它把性能副作用从"需要人看出来"变成"预算违规自动拒绝执行"，但 1\,ns 到
      1\,$\mu$s 的退化它未必能捕获，因为它计量的是抽象操作成本而非实际硬件延迟。
      实时性能保证需要时序分析与受限的运行环境，这是另一套完全不同的工程。
\end{enumerate}
\end{remark}

\section{运行时保证}\label{sec:runtime}

形式化证明依赖环境假设，而实际环境可能变化。因此除静态证明外，还需要让执行系统在
假设失效时仍然保持安全。

\begin{principle}[分层防御]\label{pr:layered-defense}
理想的架构应结合静态证明（针对已知确定性行为）与动态监控（针对未知或非确定性
行为），形成分层防御：静态证明能证明的先证明；剩余风险由独立监控检测；严重违规由
隔离或回退限制影响。
\end{principle}

\begin{remark}[监控自身的边界]\label{rmk:monitor-limit}
即使运行时监控器本身正确，也不代表它掌握了完整的执行事实。观测轨迹缺失、事件模糊、
传输不可靠与并发不可见行为都会限制结论。因此监控系统的观测机制本身必须被纳入保证链。
\end{remark}

\begin{remark}[契约组合的四个条件]\label{rmk:contract-conds}
局部性质组合依赖：环境确实满足子组件假设；连接拓扑是良置的；反馈条件可检查；
观察与时间语义兼容。组件"单独正确"与系统"整体正确"不能互换。
\end{remark}

\begin{principle}[不变量的条件式]\label{pr:ag}
一条可用的假设--保证契约应当同时写明：它依赖的环境条件、它向外保证的行为、它可以
产生的副作用、它必须消耗与占用的资源、以及它与其他子系统组合时必须满足的约束。
\end{principle}

\section{制品绑定}\label{sec:fourth}

保证条件式 \eqref{eq:assurance} 的最后一项，即实际执行的对象就是被验证的对象，在既有
讨论中经常被遗漏。它不是边角细节：即使验证器正确地接受了某个二进制，如果部署系统
使用的是另一个文件，或者加载时的环境改变了关键前提，前面的验证结论不能自动套用。
供应链攻击恰好瞄准这条裂缝。

因此完整的信任链条必须包含：内容寻址的构建产物、验证证据与二进制的密码学绑定、
以及部署时的再次校验。用式 \eqref{eq:binding} 的记号，即
$\mathrm{digest}(B_{\text{deployed}}) = \mathrm{digest}(B_{\text{verified}})$
且 $E_{\text{deployed}} \models A_{\text{claimed}}$。

\begin{principle}[可复现构建的三层区分]\label{pr:reproducible}
应当区分三种不同强度的性质：
\begin{enumerate}[label=(\alph*)]
\item \term{可复现}（reproducible）：同样的输入产生逐字节相同的输出；
\item \term{可验证}（verifiable）：独立第三方可以从输入重新构建并比对；
\item \term{已验证}（verified）：证明某个特定二进制满足其声明的性质。
\end{enumerate}
(a) 与 (b) 是 (c) 得以生效的前提条件；(c) 不含 (a)。三者不可互相替代。
\end{principle}

\section{共识层的可靠性界}\label{sec:consensus}

第 \ref{sec:trust} 节指出，意图层不存在机械解。那么它存在什么？本节给出一个精确的
答案，并证明一个反直觉的结论。

\begin{theorem}[共识可靠性界]\label{th:condorcet}
设 $n$ 个验证者对"某规格是否忠实表达意图"这一命题独立作出判断，每人判断错误的
概率至多为 $p < 1/2$。则多数表决出错的概率满足
\begin{equation}\label{eq:chernoff}
\Pr\bigl[\text{多数出错}\bigr] \;\le\; \exp\Bigl(-2n\Bigl(\tfrac{1}{2} - p\Bigr)^{2}\Bigr).
\end{equation}
\end{theorem}

\begin{proof}
设 $X$ 为出错人数，$X = \sum_{i=1}^{n} X_i$，$X_i$ 独立且
$\Pr[X_i = 1] \le p$，故 $\mathbb{E}[X] \le np$。多数出错等价于
$X \ge n/2$。由 Hoeffding 不等式\cite{hoeffding1963}，对 $t > 0$，
$\Pr[X - \mathbb{E}[X] \ge t] \le \exp(-2t^{2}/n)$。取
$t = n(1/2 - p)$（此时 $X \ge n/2$ 蕴含 $X - \mathbb{E}[X] \ge t$），得
\[
\Pr\Bigl[X \ge \tfrac{n}{2}\Bigr] \;\le\;
\exp\Bigl(-\frac{2n^{2}(1/2 - p)^{2}}{n}\Bigr) \;=\;
\exp\Bigl(-2n\bigl(\tfrac{1}{2} - p\bigr)^{2}\Bigr). \qedhere
\]
\end{proof}

这一结论是陪审团定理\cite{condorcet1785}的定量形式。但它依赖独立性，而独立性是脆弱
假设；其失效有一个同样精确的形式。

\begin{proposition}[相关错误模式支配界]\label{th:correlated}
设 $n$ 个验证者的错误具有如下结构：以概率 $q$ 发生一个\emph{共同错误模式} $C$
（例如共享的训练数据偏差、共享的领域误解、共享的利益激励），此时全部验证者同时出错；
以概率 $1-q$ 各验证者独立出错，错误率 $p' < 1/2$。则
\begin{equation}\label{eq:correlated}
\Pr\bigl[\text{多数出错}\bigr] \;\ge\; q .
\end{equation}
特别地，当 $q$ 有正下界时，增加 $n$ 不能把多数出错的概率压到 $q$ 以下。
\end{proposition}

\begin{proof}
当 $C = 1$ 时，全部 $n$ 个验证者出错，多数必然出错。故
$\Pr[\text{多数出错}] \ge \Pr[C = 1] = q$。对 $n$ 取极限不改变该下界。
\end{proof}

\begin{corollary}[多样性而非数量]\label{cor:diversity}
由命题 \ref{th:correlated}，共识层可靠性的支配变量是\emph{错误的独立性}
（即 $q$ 的大小），而不是验证者的数量。因此共识设计的核心目标是降低 $q$：
引入不同背景的参与者、不同训练来源的模型、不同利益立场的审查方，以及显式设计的
\emph{对抗性}角色。
\end{corollary}

\begin{remark}[这一结论的用途]\label{rmk:diversity-use}
推论 \ref{cor:diversity} 给出一个可直接用于工程判断的准则：当有人提议"再加几个
智能体来审查"时，正确的问题是"新增的审查者与现有的审查者在错误模式上是否独立"。
若不独立，增加数量只增加成本，不增加可靠性。这也解释了为什么"用另一个同源模型来
审查"在结构上不构成保证。
\end{remark}

\begin{remark}[共识模型是什么，不是什么]\label{rmk:consensus-status}
定理 \ref{th:condorcet}、命题 \ref{th:correlated} 与下文的设计效应近似，都是已声明的
概率模型\emph{之内}的陈述。它们不应被读作对真实审查过程可靠性的预测；模型与
实践之间的差距值得被明确说出，因为本文框架最可能被误用的地方正在这里。

三点保留。第一，这些模型假设有一个共同待判断的命题、在 $q$ 之外错误来源相互独立，
以及每人能力固定。真实的审定活动三条全都违反：什么算作忠实的规格本身就在争议之中，
审定者会通过交谈相互影响，而能力随所判断的条款而变化。关于陪审团定理的文献在很大
程度上就是在研究这些前提有多么脆弱——尤其是指出：在任何单一的概率概念下，独立性与
能力无法同时被证成 \cite{dietrich2008}——而共享的证据体是使表面的独立性失效的最自然的
机制，因为它使审定者的错误共享同一个原因 \cite{dietrich2004}。后一个结果正是原则
\ref{pr:q-decomposition} 所称信息同构的形式对应物。

第二，单一参数 $q$ 一般并不能刻画相关性。两位审定者可以在某一条款上相关、在另一条款上
独立，可以在某个子领域上因共享训练而相关、在另一个子领域上因共享激励而相关，而一维的
汇总把这一切平均掉了。$q$ 是界所需要的正确参数，同时是对该现象的有损汇总。

第三，也是后果最重的一点：除非"忠实于意图"有一个稳定的评价标准、一个被指定的裁决者，
以及一份关于什么算作错误的说明，$p$ 与 $q$ 不能从数据中估计。在这些东西存在的地方——
验收标准可观察的有界领域——这些量可以从结果中估计，这正是假设 $H_8$ 与 $H_{11}$ 所
处理的情形。在它们不存在的地方，界仍然是一个条件性陈述，而再多的审定数据也补不上
那个缺失的标准。诚实的总结是：若测量可能，模型告诉你该测什么；若测量可能，它也告诉你
答案的形状。
\end{remark}

\begin{remark}[部分相关：设计效应]\label{rmk:design-effect}
定理 \ref{th:condorcet} 与命题 \ref{th:correlated} 描述两个极端：错误独立，给出随 $n$
指数改善的界；共同错误模式，给出任何 $n$ 都无法突破的下限 $q$。真实的审定网络位于两者
之间，其错误在成对的层面上部分相关。标准统计工具在此适用，值得写下来，因为它使推论
\ref{cor:diversity} 成为定量的，而不只是修辞性的。

设 $\bar\rho$ 为审定者错误的平均成对相关。惯用的\emph{设计效应}近似 \cite{kish1965}
把名义审定者数量替换为有效数量，
\begin{equation}\label{eq:design-effect}
n_{\mathrm{eff}} \;=\; \frac{n}{1 + (n-1)\bar\rho},
\end{equation}
并以 $n_{\mathrm{eff}}$ 代替 $n$ 来读取定理 \ref{th:condorcet} 的界。由此得到两项后果。
\begin{enumerate}[label=(\alph*)]
\item 一旦 $n\bar\rho \gg 1$，式 \eqref{eq:design-effect} 就给出
      $n_{\mathrm{eff}} \to 1/\bar\rho$，\emph{与 $n$ 无关}。越过该点之后，再增加
      审定者什么也买不到；天花板只由相关性设定。
\item 降低 $\bar\rho$ 会抬高该天花板，而这一取舍是显式的：一旦 $n\bar\rho$ 很大，
      把 $\bar\rho$ 减半大致可以把可达到的 $n_{\mathrm{eff}}$ 翻倍。多样性是一条预算线，
      不是一句口号。
\end{enumerate}
它与命题 \ref{th:correlated} 的关系是：$q$ 是 $\bar\rho \to 1$ 的退化情形——完全相关时
$n_{\mathrm{eff}} \to 1$，式 \eqref{eq:design-effect} 复现了下限。

本条注附带两点保留。第一，式 \eqref{eq:design-effect} 是可交换相关下均值方差的标准近似，
而不是关于多数错误的定理；要为中间区域给出严格的界，需要一个显式的相关模型，本文不
声称给出这样的模型。第二——而这一点正是该近似可用的原因——$\bar\rho$ 可以从本文框架
已经在收集的数据中\emph{估计}出来：审定者错误的历史成对相关，以随后观察到的缺陷与
返工来结算。这并没有回答测量 $q$ 的开放问题 \ref{op:q}；它提供了一个可测量的伴随量，
而 $q$ 仍然是那个不可测量的极限。
\end{remark}

\subsection{为什么两个层无法合并，以及这强制了什么}\label{sec:enrichment}

定义 \ref{def:two-layer} 断言：这条链有一个机械层和一个共识层，而两者是被缝合起来的，
而不是被熔合在一起的。这一断言有一个数学理由，而把它说出来就解决了一个本来要靠偏好来
回答的设计问题：为什么一个能力更强的单一验证者不能取代这一对。

理由是：两个层不只是强度不同；它们把自己的保证指派给\emph{不同种类的对象}。机械保证是
一个证明：一条箭头要么带着一个证明，要么不带着，而一个证明可以精化另一个。共识保证是
一个数：一个不可靠度 $\varepsilon$，谈不上一个 $\varepsilon$ 精化另一个。在形式上，两个
层是对同一套阶段与制品结构的两种\emph{富化}——同态对象在前一种情形取自二元序，在后一种
情形取自 Lawvere quantale $V = ([0,\infty],\geq,+,0)$，其中两个阶段的复合的不可靠度至多
等于二者之和。那条复合律就是并集界，值得强调的是它不需要独立性假设：若阶段 $1$ 与 $2$
各自以至少 $1-\varepsilon_i$ 的概率忠实，那么无论二者之间有何种相依性，二者同时忠实的
概率至少是 $1-\varepsilon_1-\varepsilon_2$。这值得与定理 \ref{th:condorcet} 的界相对照：
后者的指数改善要求独立性，而第 \ref{sec:consensus} 节整节的存在就是为了对之作限定。

两种富化之间的不对称并不是对称的。存在一个从证明到可靠性的换底——已清偿的义务不可靠度
为 $0$，未清偿的义务不可靠度为 $\infty$——所以一个证明可以被正当地弱化为一个可靠性
主张。反过来不成立，而不成立的理由只有两行。若
$G : ([0,\infty],\geq,+,0) \to (\{0,1\},\leq,\wedge,1)$ 尊重幺半结构，则对每个
$\varepsilon$，
\begin{equation}\label{eq:doubling}
G(2\varepsilon) \;=\; G(\varepsilon+\varepsilon) \;=\; G(\varepsilon)\wedge G(\varepsilon)
\;=\; G(\varepsilon),
\end{equation}
所以 $G$ 对任何 $n$ 都无法区分 $\varepsilon$ 与 $2^n\varepsilon$。把这一计算弱化并不会
丢失任何东西：仅函子性就给出 $G(2\varepsilon) \leq G(\varepsilon)$，而比较映射给出反向
不等式，所以 \eqref{eq:doubling} 对松弛幺半的 $G$ 与对强的 $G$ 一样成立。配套开发 \cite{assaylean} 中所检查
的形式是严格的那一个，$G(a+b) = G(a)\wedge G(b) \Rightarrow G(2\varepsilon) =
G(\varepsilon)$；松弛的加强就是刚刚给出的两行论证，它本身并未被检查。

从 $V$ 到二元序的幺半映射恰好有三个：到顶元素的常值映射；映射
$\varepsilon \mapsto [\varepsilon < \infty]$，它是二元序包含的左伴随，也是 Lawvere 的
$\pi_0$；以及映射 $\varepsilon \mapsto [\varepsilon = 0]$，它是右伴随 \cite{lawvere1973}。
\emph{这三个在有限正值上各自都是常值}——前两个在那里取顶元素，第三个取底元素——所以它们
没有一个能分开两个这样的值。这就是本章所使用的后果：\emph{没有
幺半映射能分开两个有限正值的不可靠度}。因此对 $\varepsilon$ 施加单一阈值不只是一个有损
的汇总；它根本不是一次换底。取 $\theta$ 为阈值、$G(\varepsilon) = [\varepsilon \leq
\theta]$，则两个 $\varepsilon = 0.6\,\theta$ 的阶段各自通过，而它们在 $1.2\,\theta$ 处的
复合不通过，所以 $G$ 违反 \eqref{eq:doubling}，并随之违反那条复合律。

\begin{principle}[可靠性靠预算，而不是靠门控]\label{pr:budget}
一条必须\emph{全部}成立的 $k$ 阶段链，其复合不可靠度至多为
$\sum_{i=1}^{k}\varepsilon_i$。因此把每个阶段对固定阈值 $\theta$ 作门控，得到的复合界是
$k\theta$，它随 $k$ 无界增长；界必须改为被\emph{划分}，即在总预算 $\Theta$ 之下要求每个
阶段满足 $\varepsilon_i \leq \Theta/k$。
\end{principle}

该原则所排除的做法足够常见，值得用数字把它说出来。取"$99\%$ 可靠"作为逐条款的门控，
$\varepsilon \leq 0.01$，并取一份含 $k$ 个条款的规格，每个条款各自单独审定——是分别审定，
而不是其错误相互独立，因为并集界本就不假定独立性。复合界 $1-k\varepsilon$
在 $k = 10$ 时为 $0.90$，在 $k = 50$ 时为 $0.50$，而在 $k = 100$ 时\emph{为零}——完全
没有信息。真实的规格有数百个条款，所以按固定的逐条款可靠性作逐条款门控不只是削弱保证；
越过一百个条款之后它什么也没说，却继续看起来像一条保证。界必须被划分：要在 $k = 100$ 时
有任何内容，逐条款的要求是 $\varepsilon < 0.01$，严格低于 $1/k$。这与注
\ref{rmk:design-effect} 的预算线形状相同——那条注已经说多样性是一条预算——只是用在了
另一个量上；它也是设计效应修正取"跨阶段"而非"阶段内"的形式：设计效应管的是\emph{单个
阶段的审定者之间}的相关性，并集界管的是\emph{跨阶段}的复合。

关于这一主张有两条边界条件。第一，它关乎\emph{可靠性}的复合；它并没有说阈值无用，只是
说按阶段施加的阈值不是一个可复合的汇总，而正确的复合陈述是对每个 $\theta$ 而非某一个
选定的 $\theta$ 取 $\varepsilon$ 索引族 $[\varepsilon \leq \theta]$。该族正是定量陈述在
布尔逻辑中被表达而不丢失内容的办法，这一手法在定量等式逻辑中是标准的 \cite{bacci2023}。
单一阈值的失效模式在双模拟度量文献中也有一个标准名称：近似双模拟一般不具有传递性，而
这正是 $0.6\theta + 0.6\theta > \theta$ 这一计算的抽象形式 \cite{gebler2013}。

第二，该原则界定的是：一条链能够主张什么，而不是单个阶段能够达到什么。
$\varepsilon_i$ 是给定一种设计与一种相关结构时一轮审定的性质；设计效应则是估计它们的
方式。因此这里的贡献是狭窄的，并且应当被狭窄地陈述：不是一个更好的界，而是一个观察
——常常被当作粗糙退路的求和界，正是被富化所\emph{强制}的那条复合律；而它的替代品不是
更好的界，而是另一种对象。这一点值得说出来，因为在框架恰好需要它的那一点上，通行的
看法正好相反。Rushby 在与本章最接近的保证案例分析中，指出概率证据与演绎论证之间的
错配，给出求和规则与乘积规则，判定两者都"过于粗糙且悲观"，并把修补留作开放问题
\cite{rushby2015}。本节提供的精化是：求和规则不是一个有待改进的粗糙近似——它就是
不可靠度的复合\emph{所是}的东西——而表面的粗糙在于逐阶段的 $\varepsilon$，即设计效应
修正与第 \ref{sec:mechanism} 节的激励分析发挥作用的地方，而不在于组合它们的规则。

与软件工程文献有一个接触点值得记下，因为它是同一现象的经验对应物。取自同一利益类别的
审定者会一起失败，而对此的标准模型是设计多样性：其中独立生产的版本出现同时失效的频率
远高于独立性所允许的程度；Littlewood、Popov 与 Strigini 的综述是这些模型的标准参考
文献，而 Knight 与 Leveson 的实验是这一经验性反驳的源头
\cite{littlewood2001,knight1986}。对照 \eqref{eq:doubling} 来读，那批文献是在测量注
\ref{rmk:design-effect} 意义上的 $\bar\rho$，而不是在给它定界；这正是为什么本框架依赖
一个可估计的伴随量而不是依赖 $q$ 本身是一个务实的选择，而不只是一项让步。

\section{意图层的激励结构}\label{sec:mechanism}

命题 \ref{th:correlated} 把共识可靠性归因于参数 $q$，但没有给出 $q$ 的估计方法。
本节指出这一遗漏的性质：$q$ 的主要来源是\emph{合谋}与\emph{共同利益}，因此它不是一个
统计参数，而是一个\emph{激励设计}问题。附录 \ref{app:protocol} 把由此得出的要求转成一套
程序，并写明它自身的失效模式；那是一个**未被建造**的机制草案，不是已被验证的机制，之所以
写出来，是为了让"这个交叉地带尚未被占据"这一主张有一个具体的对象可供评判。

\begin{principle}[$q$ 的构成]\label{pr:q-decomposition}
共同错误模式可分解为三类来源：
\begin{enumerate}[label=(\alph*)]
\item \term{认知同构}：验证者共享训练来源、领域假设或思维定式；
\item \term{利益同构}：验证者的收益与"通过"相关，而非与"正确"相关；
\item \term{信息同构}：验证者读取同一份被审查材料，故其误差由同一处缺陷决定。
\end{enumerate}
三类来源对应三种不同的干预：引入异质参与者（对 (a)）、改变收益结构（对 (b)）、
引入独立信息通道（对 (c)）。只做第一类干预是常见做法，也是效用最低的做法。
第 \ref{sec:authority} 节指出同一族中的第四种共现，\emph{权威同构}，并解释它为什么不是
本分解的第四项，而是一条独立的失效轴。
\end{principle}

\begin{principle}[适当评分规则]\label{pr:scoring}
对第二类来源，正解是让验证者的报酬与其判断的\emph{准确性}而非其判断的\emph{倾向}
挂钩。适当评分规则\cite{brier1950}提供了现成工具：报酬函数在真实信念处取唯一极大值，从而对
"迎合多数"与"讨好提交方"两种行为同时去激励。具体形式可以是一个针对"该规格是否忠实
表达意图"的预测市场，其结算依据是后续观察到的缺陷与返工数据。该机制有两个副产品：
\begin{enumerate}[label=(\alph*)]
\item 它为降低 $q$ 提供了除"换人"之外的手段；
\item 它给出 $q$ 的后验估计，测量验证者误差的历史相关性即可。
\end{enumerate}
\end{principle}

这一节补上了第 \ref{sec:consensus} 节留下的缺口。那里的界是条件性的，其条件
（独立性）在此被归约为一个可设计、可测量、可付费的量。第 \ref{ch:program} 章把
它写成假设 $H_{11}$。

\section{谁可以审定：道义轴}\label{sec:authority}

第 \ref{sec:consensus} 节与第 \ref{sec:mechanism} 节把审定当作一个认识论问题：规格是否
忠实于意图，以及说"是"的多数有多可靠。还存在第二个问题，而本文框架自己的保证条件式
已经包含了它，只是没有把它接上。式 \eqref{eq:assurance} 的第二个合取项要求\emph{契约
已明确且被授权采用}。授权不是可靠性问题：一份规格可以被一个规模庞大且真正独立的专家
组忠实地审定，同时仍然是无人具备采用资格的那一份。

\begin{principle}[授权是一条独立的失效轴]\label{pr:authority}
审定网络的可靠性关乎其裁决是否\emph{很可能是正确的}。授权关乎该裁决是否\emph{决定得了
任何事情}。两者相互独立。
\begin{enumerate}[label=(\alph*)]
\item 一个裁决可以既可靠又未被授权：一个完全来自同一利益相关者类别的、有能力的专家组，
      就另一类别有资格拒绝的一条需求达成一致。
\item 一个裁决也可以既被授权又不可靠：单一责任主体在信息不足的情况下作出判断。
\end{enumerate}
因此，$q$ 的改善、$n_{\mathrm{eff}}$ 的改善或激励结构的改善，都无法修复一次授权失效。
这种失效不是所有人一起错了，而是所有人加在一起并不是有权作出决定的那个机构。
\end{principle}

这就是权威应当位于 $q$ 的三类来源\emph{之旁}、而不是其中的原因。原则
\ref{pr:q-decomposition} 把共同错误模式分解为认知同构、利益同构与信息同构——审定者以
同一种方式出错的三种途径。第四种共现很容易命名，称为\term{权威同构}：全部审定者都来自
同一个利益相关者类别。它确实会抬高 $q$，因为同一个类别往往既共享前提又共享利益。但
降低 $q$ 并不解决它。一个完全多样化的外部人专家组，恰恰是授权失效最干净的情形：
他们可能是对的，却仍然没有资格。

这一区分有一种用旧了的制度形态，值得用一句话、也只值得一句话。分布式账本系统把确立
\emph{什么是真的}的机制，与确立\emph{谁可以决定什么算作真的}的机制分开；共识与治理是
不同的协议，一个共识极佳却没有治理的系统并不因此就是正当的。这个类比只是作为类比给出
的。它指向的是本文的一个缺口，而不是本文的一个结果。

\begin{remark}[对意图层设计的两项后果]\label{rmk:authority-design}
\begin{enumerate}[label=(\alph*)]
\item 第 \ref{sec:consensus} 节的审定记录，应在来源之外同时携带每位审定者的\emph{利益
      相关者类别}。这多花一个字段，使权威同构变得可观察，并使授权失效在事后可以与
      可靠性失效区分开来——两者的补救手段不同。
\item 被授权机构的选择不能从本文框架中推出。它是关于"谁承担决定的后果"的制度事实，
      并且由与第 \ref{ch:levels} 章把 $\Lk{\infty}$ 指派给人来承担的那同一个论证所
      确定——而那个论证的前提是现行的法律责任制度，正如该章明确指出的。
\end{enumerate}
\end{remark}

\section{规格绑定}\label{sec:spec-binding}

式 \eqref{eq:assurance} 的第四支柱把已验证的制品绑定到实际执行的制品。到此为止还没有
任何东西把规格绑定到它本应表达的那个意图，而这种不对称是要紧的，因为两者随时间的行为
相反。一个二进制要么逐字节相同、要么不相同：它的绑定是静态的，并且一次判定即定。而
规格所要忠实的是一个会移动的意图。命题 \ref{prop:three-grounds}(i) 就是原因——需求是
生成性的，偏好由规格写定时尚不存在的情境激活——所以忠实性不是规格所拥有的一种性质，
而是一个会衰减的量。

表 \ref{tab:evidence-package} 已经要求每个制品携带第 10 项，即重新验证规则：哪些变化
触发重新计算受影响的证据。规格一侧没有对应的东西，而命题 \ref{prop:assurance} 的条件
(i) 被陈述得好像它不需要任何对应物。修法是让规格自己携带它的绑定记录。

\begin{definition}[规格绑定]\label{def:spec-binding}
一个\term{规格绑定}是一个四元组
\[
\mathrm{SB}(\Phi) \;=\; \bigl(\text{规格版本},\; \text{审定记录},\;
\text{有效视界},\; \text{重新审定触发条件}\bigr),
\]
其中审定记录是第 \ref{sec:consensus} 节的共识制品，它在规格写定时确立了 $\Phi$ 的
忠实性；有效视界是该审定被当作仍然现行的那个区间；重新审定触发条件则是在视界到期之前
使其失效的那些事件。
\end{definition}

触发条件才是有意思的部分，而且它们不是随意的。命题 \ref{prop:three-grounds} 命名了
三种漂移来源，每一种都提示出自己的触发条件。
\begin{enumerate}[label=(\alph*)]
\item \emph{生成性需求}：出现了新的一类用户、场景或部署环境，而这恰恰是激活了无人曾
      表达过的偏好的那种情境。
\item \emph{有损压缩}：写 $\Phi$ 所用的领域模型、本体或受控语言发生变化，以致同一段
      文字所指的东西变了。
\item \emph{确认}：一起被观察到的事故，其处置无法从 $\Phi$ 推出，这是 $\Phi$ 对意图
      表达不足、而不是实现出了错的证据。
\end{enumerate}
触发条件 (c) 值得单独指出，因为它是免费得到的：一起需要人来作出判断的事故，按构造
就是规格没有解决的事故。

因此，命题 \ref{prop:assurance} 的条件 (i) 应当带着一个时间下标来读。记
$\varepsilon_{\text{cons}}(\Delta t)$ 为自审定起经过区间 $\Delta t$ 之后的共识错误，
\begin{equation}\label{eq:spec-decay}
\Pr\bigl[\Phi \text{ 在时刻 } t \text{ 忠实于 } I\bigr] \;\ge\;
1 - \varepsilon_{\text{cons}}(\Delta t),
\qquad \Delta t = t - t_{\text{attest}},
\end{equation}
它在 $\Delta t = 0$ 时退化为原来的条件。本文不给出 $\varepsilon_{\text{cons}}(\cdot)$
的形式：不提出任何衰减模型，该量被留作一个可测量的对象，而不是被建模的对象。本节所
确立的是结构上的对称性——规格需要一条重新审定规则，正如二进制需要一条重新验证规则——
以及这一对称性是由命题 \ref{prop:three-grounds}(i) 逼出来的，而不是为了完备而添上的。

\begin{remark}[制品绑定缺口的规格一侧对偶]\label{rmk:spec-h9}
假设 $H_9$ 关乎制品绑定缺口的可测量风险。它的规格一侧对偶是"依据一份陈旧但仍被审定的
规格执行"的风险：审定记录仍在有效视界之内，二进制仍与其摘要相符，而两者一起错了，
因为意图移动了。这两个缺口相互独立，一次部署可以堵住其中一个而让另一个敞着。
\end{remark}

\section{综合保证条件}\label{sec:assurance-final}

把本章各部分组合起来，得到一个完整的、条件性的保证陈述。

\begin{definition}[双层信任结构]\label{def:two-layer}
从意图到执行的链条是一个双层信任结构的复合体：
\begin{equation}\label{eq:twolayer}
\underbrace{I \xrightarrow{\;\text{共识}\;} \Phi}_{\text{不可形式化层}}
\;\longrightarrow\;
\underbrace{\Phi \xrightarrow{\;\text{机械}\;} P \xrightarrow{\;\text{机械}\;} B
\xrightarrow{\;\text{机械}+\text{监控}\;} E}_{\text{可形式化层}}
\end{equation}
在可形式化层，信任由机械验证器保证；在不可形式化层，信任由对抗性共识保证；
两层之间的边界、共识的记录以及整个链条的可追溯性，由密码学锚定的基础设施机械保证。
\end{definition}

\begin{proposition}[综合保证：充分条件]\label{prop:assurance}
设意图 $I$、已审定规格 $\Phi$、候选实现（二进制）$B$、运行环境 $E$，
$A$ 为声明的环境假设。若下列条件同时成立，则"在所声明的范围 $\Sigma$ 内满足性质
$\varphi$"成立：
\begin{align}
&\text{(i)}\quad \Pr\bigl[\Phi \text{ 在时刻 } t \text{ 忠实于 } I\bigr] \;\ge\;
1 - \varepsilon_{\text{cons}}(\Delta t);\\
&\text{(ii)}\quad \Tr(B, E) \;\subseteq\; \Allow(\Phi);\\
&\text{(iii)}\quad E \models A;\\
&\text{(iv)}\quad \mathrm{digest}(B_{\text{deployed}}) = \mathrm{digest}(B_{\text{verified}});\\
&\text{(v)}\quad \varphi \text{ 在所用的解释关系 } R \text{ 下被保持};\\
&\text{(vi)}\quad \Phi \text{ 由一个具备采用资格的机构授权},
\end{align}
其中 $\varepsilon_{\text{cons}}(\cdot)$ 是式 \eqref{eq:chernoff} 的界，由 $q$ 与 $\bar\rho$
决定，并在 $\Delta t = t - t_{\text{attest}}$ 处取值。
\end{proposition}

这一陈述是刻意单向的。它在别处曾被写成等价式，而那是错的；其理由值得说出来，因为
只要提出这种形状的框架，这个理由就会重现。这六个条件是\emph{证据}条件：它们说的是
手上必须有什么。它们不是真值条件，而真值并不等待它们。一条性质可以在没有任何证据
被汇集起来的情况下成立，所以反方向之所以失败，理由是最基本的那一条：一个系统可以是
正确的却没有留下文档。本命题所断言的是有用的那个方向——六条全部成立就足够了——而它
刻意不断言的是：少于六条就不可能。

每个合取项都带有自己的保留，而它们没有一个与其余各项在相同的意义上属于机械性的。
\begin{enumerate}[label=(\alph*)]
\item 条件 (i) 是一个概率，在忠实性被定义、裁决者被指定、错误类别被说明之前，它没有
      自然的定义。在这些东西存在之前，$\varepsilon$ 是模型的参数，而不是测量到的
      频率。它的可靠性受 $q$ 以及第 \ref{sec:consensus} 节与第 \ref{sec:mechanism} 节的
      相关 $\bar\rho$ 支配，而两者都是条件模型的量。
\item 条件 (ii) 是一条语义性质，不是例行检查。$\Tr(B,E) \subseteq \Allow(\Phi)$ 只有在
      受限语言、有限模型或携带可检查证书的程序上才是直接可判定的；对一般环境中的一般
      程序，它是一项证明义务，靠类型化、模型检查、契约验证与不同强度的测试组合来
      清偿——前提是它真被清偿了。
\item 条件 (iii) 与 (iv) 是那对机械性的条件，其强度受可信计算基与环境模型限制，而不是
      受本章中任何东西限制。
\item 条件 (v) 正是注 \ref{rmk:relational-compose} 的关系性前提被花掉的地方。
\item 条件 (vi) 是道义性的，根本不是可靠性问题（第 \ref{sec:authority} 节）。
\end{enumerate}
因此，定义 \ref{def:assurance} 的六个合取项中，三个是机械性的，一个是概率性的，一个是
道义性的，而还有一个——契约必须明确这一要求——正是本章其余部分的用途所在。

综合陈述还必须携带另一个区分，即经典的安全性（safety）与活性（liveness）之分。条件
(ii) 约束的是有限轨迹：被禁止的事情不会发生。进展性质——一个请求最终得到应答、一把锁
最终被释放、一个队列最终排空——是活性性质，它们不是"没有有限前缀违反 $\Phi$"的实例，
也不能由同一个论证得到。当系统的契约包含活性时，这些义务属于能够表达它们的层级与
装置：对协议模型是公平性假设与良基性论证，对已部署的系统则是带界视界的监控。一份把
活性悄悄吸收进轨迹包含的综合陈述，会以一种容易被忽略的方式出错，而本陈述没有这样做。

式 \eqref{eq:twolayer} 与命题 \ref{prop:assurance} 一起给出了对问题 \ref{prob:trust}
的完整回答：

\begin{principle}[对问题 \ref{prob:trust} 的回答]\label{pr:answer}
"人不再读代码"是可达的；"人不必再为目标的表达负责"是不可达的。可行的分工是：
人确立并承担目标与约束的责任；生成器负责大规模生成、搜索与优化；机械化的验证与
执行边界负责判断与限制实现；系统始终携带可追溯的保证证据。
\end{principle}

\begin{remark}[中心问题 Q1--Q8 的结论摘要]\label{rmk:q1q8}
\begin{longtable}{L{5.4cm}L{9.0cm}}
\toprule
\textbf{问题} & \textbf{本文的判断} \\
\midrule
\endhead
Q1 能否证明意图与规格一致？ &
  不能依靠通用自动方法。受控语言、领域本体、例子与人工审定可缩小问题范围，
  但无法覆盖未表达的隐含需求（命题 \ref{prop:three-grounds}）\\
Q2 如何证明规格充分？ &
  可检查一致性、可满足性、可实现性、空洞成立与形式化覆盖指标；这些都不能证明规格
  包含了所有未声明的禁令 \\
Q3 局部保证如何推出端到端保证？ &
  需要关系可复合、环境假设闭合、观察映射兼容（命题 \ref{prop:soundness}）；
  活性与实时性质通常还需额外条件 \\
Q4 能否统一功能与资源保证？ &
  可以在联合行为--成本语义中表达多种约束，但时间、内存、能耗与信息流需要不同的
  分析模型；联合表达不等于统一求解 \\
Q5 如何降低对生成模型的信任？ &
  生成器提出候选，独立检查器验证。这降低了对生成器正确性的依赖，但不消除对逻辑
  内核、定理陈述、外部假设与制品对应关系的信任要求 \\
Q6 给定规格是否总能有效综合？ &
  不能。有限语法、受限逻辑与特定模板可使部分问题可解；任意程序语义性质存在
  不可判定性与搜索空间爆炸的边界 \\
Q7 实际延迟能否严格保证？ &
  在特定硬件、输入、调度与成本模型假设下可得到可用的最坏执行时间上界；基准测试
  与渐近复杂度本身不足以证明纳秒级上界 \\
Q8 是否需要统一理论？ &
  更有根据的选择是建立共享的证明接口与证据格式，同时保留各领域不同的语义、假设与
  分析工具。统一形式体系可以是目标之一，但不是必要的先验条件 \\
\bottomrule
\end{longtable}
\end{remark}

本章的经验内容是方向 1--3（第 \ref{ch:program} 章）与假设 $H_7$、$H_8$、$H_{11}$。
命题 \ref{prop:no-godel} 与 \ref{prop:three-grounds} 是论证而非定理，其可挑战之处
在于三条根据是否穷尽。定理 \ref{th:condorcet} 与命题 \ref{th:correlated} 是标准
概率论结果的应用，其结论依赖独立性假设。
